\documentclass[11pt]{article}
\usepackage[T1]{fontenc}
\usepackage[a4paper,margin=1in]{geometry}
\usepackage{amsmath,amssymb,amsthm}
\usepackage{graphicx}
\usepackage[hidelinks]{hyperref}
\usepackage{microtype}

\newtheorem{theorem}{Theorem}
\newtheorem{lemma}[theorem]{Lemma}
\newtheorem*{sourcequestion}{Source question}

\newcommand{\Z}{\mathbb Z}
\newcommand{\N}{\mathbb N}
\newcommand{\E}{\mathbb E}

\setlength{\parindent}{0pt}
\setlength{\parskip}{0.55em}

\title{Positive-Entropy Systems Admit Divergent\\
Non-Commuting Double Averages}
\author{Candidate partial solution packet\\
Run \texttt{fa\_banach\_001} --- lane 0}
\date{9 August 2026}

\begin{document}
\maketitle

\begin{center}
\fbox{\parbox{0.9\textwidth}{
\textbf{Status: candidate partial result, likely valid.}
Austin's question has an affirmative answer for every weakly mixing
probability-preserving automorphism of positive entropy.  More generally, the
proof needs only ergodicity and positive entropy.  A subsequent identification
shows that the universal affirmative question is false: Chacon's weakly mixing
transformation forces convergence against every second system.  The theorem
here remains a positive-entropy occurrence result.
}}
\end{center}

\section{Source problem}

Let $S$ and $T$ be invertible probability-preserving transformations of a
standard probability space $(X,\mathcal B,\mu)$, and let $f\in L^\infty(\mu)$.
Austin~\cite{AustinDivergence} studies the scalar averages
\begin{equation}\label{eq:averages}
 A_N(S,T;f):=\frac1N\sum_{i=0}^{N-1}
 \int_X f(S^i x)f(T^i x)\,d\mu(x).
\end{equation}
His examples include isomorphic rigid Gaussian transformations $S,T$ for
which these averages diverge.  The conclusion of the paper asks the following
broader question (arXiv PDF page~6).

\begin{center}
\includegraphics[width=0.98\textwidth]{figures/open_problem_crop.png}
\end{center}

\begin{sourcequestion}
Can every weakly mixing transformation appear, up to isomorphism, as $S$ in a
non-convergent example of the averages~\eqref{eq:averages}?
\end{sourcequestion}

The theorem below settles every positive-entropy instance.  It says nothing
about zero entropy; a separate literature-implied obstruction now shows that
no uniform affirmative extension to all weakly mixing systems is possible.

\section{Proof intuition}

The weak Pinsker theorem~\cite{AustinPinsker} decomposes every
positive-entropy ergodic system into a direct product of a nontrivial Bernoulli
shift and a remainder.  The only new construction needed is therefore on an
arbitrary Bernoulli shift.

On a two-sided Bernoulli shift, a permutation $p$ of the coordinate labels
induces a measure-preserving coordinate permutation $C_p$.  The conjugate
$\tau=C_p^{-1}\sigma C_p$ is another copy of the same Bernoulli shift.  We
choose $p$ so that its fixed points among $0,1,\ldots,N-1$ have no asymptotic
density.  If $f$ is a centered observable of coordinate zero, then at time $i$
the two copies of $f$ inspect coordinates $i$ and $p(i)$.  Their correlation is
the variance of $f$ when those coordinates coincide and zero when they are
distinct.  Thus the correlation averages are exactly a constant multiple of
the non-convergent fixed-point densities.  No approximation is involved.

\section{The theorem}

\begin{theorem}\label{thm:main}
Let $(X,\mathcal B,\mu,S)$ be an ergodic invertible
probability-preserving system with positive Kolmogorov--Sinai entropy
$h_\mu(S)>0$, allowing $h_\mu(S)=\infty$.  Then there are an invertible
probability-preserving transformation $T$ of $X$ and a bounded real function
$f\in L^\infty(\mu)$ such that:
\begin{enumerate}
\item $T$ is measurably isomorphic to $S$; in fact, $T$ is conjugate to $S$ on
the same probability space;
\item the sequence $A_N(S,T;f)$ in~\eqref{eq:averages} does not converge.
\end{enumerate}
More precisely, $f$ may be chosen so that
\[
 \liminf_{N\to\infty}A_N(S,T;f)=0,
 \qquad
 \limsup_{N\to\infty}A_N(S,T;f)=\|f\|_2^2>0.
\]
Consequently, Austin's question has an affirmative answer for every weakly
mixing transformation of positive entropy.
\end{theorem}

We first isolate the elementary Bernoulli construction.

\begin{lemma}\label{lem:permutation}
There is a permutation $p:\Z\to\Z$ with $p(0)=0$ for which
\[
 d_N:=\frac1N\#\{0\leq i<N:p(i)=i\}
\]
satisfies $\liminf_N d_N=0$ and $\limsup_N d_N=1$.
\end{lemma}

\begin{proof}
Set $N_0=1$.  Recursively choose an even integer $L_k>kN_k$ and put
$N_{k+1}=N_k+L_k$.  Thus every block
$I_k=[N_k,N_{k+1})\cap\Z$ has even cardinality.  On blocks with even $k$, let
$p$ be the identity.  On blocks with odd $k$, interchange adjacent pairs, so
$p(N_k+2j)=N_k+2j+1$ and $p(N_k+2j+1)=N_k+2j$.  Finally, let $p(n)=n$ for
$n\leq0$.

These disjoint blockwise prescriptions define a permutation of $\Z$ fixing
zero.  At the endpoint of an identity block, the newest block alone occupies
more than $k/(k+1)$ of all preceding nonnegative indices, so the fixed-point
density tends to one along those endpoints.  At the endpoint of a swapped
block, all fixed points lie before that block, and hence their density is less
than $1/(k+1)$.  The two asserted subsequential limits follow.
\end{proof}

\begin{lemma}\label{lem:bernoulli}
Let $(K,\kappa)$ be any nontrivial standard probability space, let
$Y=K^\Z$ with product measure $\nu=\kappa^\Z$, and let
\[
 (\sigma y)_n=y_{n+1}\qquad(n\in\Z)
\]
be the two-sided Bernoulli shift.  There are a transformation $\tau$ conjugate
to $\sigma$ and a bounded real $g\in L^\infty(\nu)$ for which
\[
 \liminf_N A_N(\sigma,\tau;g)=0,
 \qquad
 \limsup_N A_N(\sigma,\tau;g)=\|g\|_2^2>0.
\]
\end{lemma}

\begin{proof}
Choose a nonzero bounded real function $h$ on $K$ with
$\int_Kh\,d\kappa=0$; such a function exists because the base is nontrivial.
Let $p$ be the permutation from Lemma~\ref{lem:permutation}, and define the
coordinate permutation
\[
 (C_p y)_n:=y_{p(n)}.
\]
Product measure is invariant under coordinate permutations, so $C_p$ is an
invertible measure-preserving map.  Put
\[
 \tau:=C_p^{-1}\sigma C_p,
 \qquad g(y):=h(y_0).
\]
Then $\tau$ is conjugate to $\sigma$.  Since $p(0)=0$, for every $i\geq0$,
\[
 (\tau^i y)_0
 = (C_p^{-1}\sigma^iC_p y)_0
 = y_{p(i)}.
\]
It follows that
\[
 g(\sigma^i y)=h(y_i),
 \qquad
 g(\tau^i y)=h(y_{p(i)}).
\]
If $p(i)=i$, their integral product is $\|h\|_{L^2(\kappa)}^2$.  If
$p(i)\neq i$, the two coordinates are independent and $h$ is centered, so the
integral product is zero.  Therefore
\[
 \int_Yg(\sigma^iy)g(\tau^iy)\,d\nu(y)
 =\|h\|_2^2\mathbf 1_{\{p(i)=i\}},
\]
and hence
\[
 A_N(\sigma,\tau;g)=\|h\|_2^2d_N.
\]
Lemma~\ref{lem:permutation} completes the proof.
\end{proof}

\begin{proof}[Proof of Theorem~\ref{thm:main}]
By the weak Pinsker theorem~\cite{AustinPinsker}, for every $\varepsilon>0$
the system $(X,\mu,S)$ is measurably isomorphic to a direct product
\begin{equation}\label{eq:pinsker}
 (Y\times Z,\nu\times\eta,\sigma\times R),
\end{equation}
where $\sigma$ is a Bernoulli shift and $h_\eta(R)<\varepsilon$.

If $0<h_\mu(S)<\infty$, choose $\varepsilon<h_\mu(S)$.  Entropy is invariant
under isomorphism and additive under direct products, so
\[
 h_\nu(\sigma)=h_\mu(S)-h_\eta(R)>0.
\]
If $h_\mu(S)=\infty$, the finite bound on $h_\eta(R)$ forces
$h_\nu(\sigma)=\infty$.  Thus the Bernoulli factor is nontrivial in either
case.

Apply Lemma~\ref{lem:bernoulli} to obtain $\tau$ conjugate to $\sigma$ and a
bounded real $g$ with divergent scalar averages.  On the product space set
\[
 S_0:=\sigma\times R,
 \qquad T_0:=\tau\times R,
 \qquad f_0(y,z):=g(y).
\]
If $\tau=C_p^{-1}\sigma C_p$, then
\[
 T_0=(C_p\times\mathrm{id}_Z)^{-1}S_0
     (C_p\times\mathrm{id}_Z),
\]
so $T_0$ and $S_0$ are conjugate.  Since $f_0$ ignores the $Z$ coordinate,
Fubini's theorem gives
\[
 A_N(S_0,T_0;f_0)=A_N(\sigma,\tau;g)
\]
for every $N$.  The latter sequence has liminf zero and limsup
$\|g\|_2^2=\|f_0\|_2^2$.

Finally, transport $T_0$ and $f_0$ back to $X$ through the isomorphism
in~\eqref{eq:pinsker}.  Conjugacy, boundedness, and all the displayed
integrals are preserved.  This gives the required $T$ and $f$.
\end{proof}

\section{Verification, limitations, and novelty status}

The proof is exact after the weak Pinsker splitting.  The two calculations a
reviewer should check most closely are the coordinate identity
$(\tau^iy)_0=y_{p(i)}$ and the entropy argument ensuring that the Bernoulli
factor is nontrivial.  Both are expanded explicitly above.  The accompanying
\texttt{VERIFICATION.md} gives a separate audit.

This result is partial: weakly mixing zero-entropy automorphisms include many
of the difficult rigid, rank-one, and distal-adjacent regimes motivating the
source question, and the argument says nothing about them.  A subsequent
lane-0 packet records the decisive negative obstruction: Lesigne--Rittaud--de
La Rue proved that Chacon's weakly mixing transformation is universal for weak
disjointness, which forces convergence on the same-space diagonal against
every second system~\cite[Proposition~4.2]{LRR}.  Thus the theorem here is a broad positive occurrence
result, but it cannot extend to every weakly mixing transformation.

A bounded novelty search through 9 August 2026 checked the run registry and
cheap local indexes, the exact wording of Austin's question, title/citation
matches for arXiv:2407.08630, and recent related arXiv papers 2407.13741,
2410.11787, and 2411.02024.  Those searches found Gaussian,
Poisson-suspension, simple-spectrum, and rank-one-base constructions, but no
statement of the positive-entropy reduction proved here.  This is not an
exhaustive literature review, so novelty confidence is moderate.  Human
review is recommended before treating the result as established literature.

\begin{thebibliography}{9}
\bibitem{AustinDivergence}
T. Austin,
\emph{Non-convergence of some non-commuting double ergodic averages},
Proc. Amer. Math. Soc. \textbf{153} (2025), 1701--1707;
arXiv:2407.08630.

\bibitem{AustinPinsker}
T. Austin,
\emph{Measure concentration and the weak Pinsker property},
Publ. Math. Inst. Hautes \'{E}tudes Sci. \textbf{128} (2018), 1--119;
arXiv:1705.00302.

\bibitem{LRR}
E. Lesigne, B. Rittaud, and T. de La Rue,
\emph{Weak disjointness of measure preserving dynamical systems},
Ergodic Theory Dynam. Systems \textbf{23} (2003), 1173--1198;
doi:10.1017/S0143385702001505.
\end{thebibliography}

\end{document}
